\chapter{Organização do Texto}
\label{cap:Organizacao_do_Texto}

\textcolor{gray}{O cuidado em identificar pontos críticos no acompanhamento
das preferências de consumo assume importantes posições no estabelecimento
do sistema de participação geral. Por outro lado, a expansão dos mercados
mundiais facilita a criação das diretrizes de desenvolvimento para o futuro.}

\section{Isso é Uma Seção}
\label{sec:Isso_e_Uma_Secao}

\textcolor{gray}{Por conseguinte, a percepção das dificuldades facilita a
criação do levantamento das variáveis envolvidas. Nunca é demais lembrar
o peso e o significado destes problemas, uma vez que a expansão dos mercados
mundiais representa uma abertura para a melhoria dos modos de operação
convencionais.}

\subsection{Isso é Uma Subseção}
\label{sec:Isso_e_Uma_Subsecao}

\textcolor{gray}{Não obstante, o desenvolvimento contínuo de distintas formas
de atuação exige a precisão e a definição das direções preferenciais no sentido
do progresso. Ainda assim, existem dúvidas a respeito de como a constante
divulgação das informações promove a alavancagem dos conhecimentos estratégicos
para atingir a excelência.}

%------------------------------------------------------------------------------%
\subsubsection{Isso é Uma Subsubseção}
\label{sub:Isso_e_Uma_Subsubsecao}

\textcolor{gray}{Caros amigos, o entendimento das metas propostas estimula
a padronização do impacto na agilidade decisória. Desta maneira, a hegemonia
do ambiente político apresenta tendências no sentido de aprovar a manutenção
de todos os recursos funcionais envolvidos.}

\section{Citando as Referências Bibliográficas}
\label{sec:Citando_as_Referencias_Bibliograficas}

O arquivo \texttt{refs.bib}
contém informações bibliográficas para alguns trabalhos que
costumam ser citados em dissertações do Profmat,
tais como os livros da Coleção Profmat da SBM, algumas leis e o GeoGebra.
O modelo adotado pelo Profmat no Cefet-MG é o numérico, onde
existem várias formas de fazer as citações das referências bibliográficas.
\begin{enumerate}[label=\alph*)]
  \item Citando numericamente, como no caso do livro sobre redação
        Matemática~\cite{burden2010numerical}.
  \item Citando apenas o nome do autor~\citeauthor{burden2010numerical}.
  \item Citando o nome do autor e o número~\textcite{burden2010numerical}.
  \item Citando apenas o ano do trabalho~\citeyear{burden2010numerical}.
  \item Vários trabalhos podem ser citados
        simultaneamente~\textcite{burden2010numerical,SOUZA:Bayesiana,chapra2015applied}.
  \item Citando com o número da página~\cite[10]{burden2010numerical}
        ou~\textcite[10]{burden2010numerical}
\end{enumerate}

\textcolor{red}{\textsc{Atenção}}: Essa classe utiliza o bib\LaTeX{} com o
\texttt{biber} para gerar as referências bibliográficas. Dependendo do
programa que você utilizar para editar e compilar a dissertação ele vai
identificar isso automaticamente e nada precisa ser feito.
Porém, em outros programas é necessário informar que a ferramenta
para gerar a bibliografia é \texttt{biber}. No
\href{https://www.texstudio.org}{TeXstudio} você deve selecionar
\texttt{Opções -- Configurar TeXstudio} no menu. No quadro de diálogo
que será aberto você deve selecionar a aba \texttt{Compilar}
e dentro dela alterar a opção \texttt{Ferramenta padrão de bibliografia}
substituindo \texttt{BibTeX} por \texttt{Biber}.

%------------------------------------------------------------------------------%
\section{Citações Diretas com Um ou Mais Parágrafos}
\label{sec:Citacoes_Diretas_com_Um_ou_Mais_Paragrafos}

\textcolor{gray}{O incentivo ao avanço tecnológico, assim como a complexidade
dos estudos efetuados maximiza as possibilidades por conta da gestão inovadora}
\textcite[10]{burden2010numerical}.
\begin{quote}
  \textcolor{gray}{O cuidado em identificar pontos críticos na mobilidade
  dos capitais internacionais possibilita uma melhor visão global dos modos
  de operação convencionais.}
\end{quote}
\textcolor{gray}{Assim mesmo, a constante divulgação das informações nos
obriga à análise do sistema de participação geral. Neste sentido, o novo
modelo estrutural aqui preconizado garante a contribuição}
\textcite[10]{burden2010numerical}.
\begin{quotation}
  \textcolor{gray}{Neste sentido, o novo modelo estrutural aqui preconizado
  garante a contribuição de um grupo importante na determinação das novas
  proposições.}

  \textcolor{gray}{A prática cotidiana prova que o desenvolvimento contínuo
  de distintas formas de atuação apresenta tendências no sentido da manutenção.}

  \textcolor{gray}{Percebemos, cada vez mais, que o desafiador cenário
  globalizado estimula a padronização dos índices pretendidos.}
\end{quotation}

\textcolor{gray}{A prática cotidiana prova que a crescente influência da
mídia desafia a capacidade de equalização das regras de conduta normativas.}

%------------------------------------------------------------------------------%
\section{Miscelâneas}
\label{sec:Miscelaneas}

Colocando texto entre `aspas simples' e ``aspas duplas''.

\textbf{Destacando um trecho do texto com negrito.}

\textit{Destacando um trecho do texto com itálico.}

\texttt{Destacando um trecho do texto com fonte de máquina de escrever.}

\textsc{Destacando um trecho do texto com fonte com Small Caps.}

Escrevendo números ordinais em modo texto
1\textsuperscript{o},
2\textsuperscript{o},
3\textsuperscript{o}.

Escrevendo porcentagens 25\%
e outros símbolos
\#
\$
\%
\&
\_
\{
\}
\textbackslash
\textasciitilde
\textasciicircum.

Escrevendo o endereço de uma página na internet
\url{https://pt.wikipedia.org}.

Criando uma nota de rodapé\footnote{Uma nota de rodapé.}.

\textcolor{gray}{O cuidado em identificar pontos críticos no acompanhamento
das preferências de consumo assume importantes posições no estabelecimento
do sistema de participação geral.}

%------------------------------------------------------------------------------%
%------------------------------------------------------------------------------%
%------------------------------------------------------------------------------%
\chapter{Objetos Flutuantes}
\label{cap:Objetos_Flutuantes}

Objetos flutuantes são figuras e tabelas cuja posição no texto é decidida
pelo \LaTeX{} e portanto são referenciados apenas por seus números.

%------------------------------------------------------------------------------%
\section{Figuras}
\label{sec:Figuras}

\textcolor{gray}{Por outro lado, o surgimento do comércio virtual causa impacto
indireto na reavaliação do sistema de participação geral.
O cuidado em identificar pontos críticos no desafiador cenário globalizado pode
nos levar a considerar a reestruturação dos procedimentos normalmente adotados.}

% Incluindo uma figura
\begin{figure}
  \centering
  \includegraphics[width=0.5\linewidth]{pitagoras}
  \caption[Teorema de Pitágoras]{O Teorema de Pitágoras relaciona os
  comprimentos dos lados de um triângulo retângulo.}
  \label{fig:pitagoras}
\end{figure}

\textcolor{gray}{É importante questionar o quanto a hegemonia do ambiente
político oferece uma interessante oportunidade para verificação das direções
preferenciais no sentido do progresso.}

%------------------------------------------------------------------------------%
%------------------------------------------------------------------------------%
%------------------------------------------------------------------------------%